% 使用说明：
% 1. 在主文档中通过 % 使用说明：
% 1. 在主文档中通过 % 使用说明：
% 1. 在主文档中通过 % 使用说明：
% 1. 在主文档中通过 \input{chapters/related_work_gap123} 引入本文件。
% 2. 参考文献数据库建议加入 references/related_work_gap123.bib。

\chapter{相关研究综述}

侧方泊车是低速自动驾驶中的典型任务。该任务对路径规划提出了较高要求。车辆既要在有限空间内完成无碰撞机动，又要保证轨迹平滑、转向可执行和末端停放精度。传统搜索算法易于生成可行路径，但所得结果常呈折线形态，难以直接满足车辆连续转向需求。贝塞尔曲线具有良好的平滑性与参数化表达能力，因而被广泛用于车辆轨迹优化。围绕侧方泊车中的贝塞尔曲线优化问题，梳理现有研究进展，有助于明确相关方法的演化脉络，并为后续问题建模与算法设计提供依据。

\subsection{侧方泊车相关研究现状}

现有泊车路径规划研究，主要面向自动代客泊车与停车场低速机动场景\cite{yin2025avp,wu2025agvparking,zhao2024automatedparkingbev}。相关工作多以环境建模、初始路径搜索和轨迹跟踪控制为主线展开。研究目标已从静态最短路搜索，转向考虑车辆运动学、障碍分布和控制执行约束的综合规划。其中，尹禹兮围绕自主代客泊车构建了搜索与优化结合的分层规划思路，吴宇轩则以智能停车场 AGV 为对象，讨论了路径规划与路径跟踪的一体化实现，Zhao 进一步把视觉 BEV 局部建图引入自动泊车规划过程\cite{yin2025avp,wu2025agvparking,zhao2024automatedparkingbev}。

从方法演化看，泊车研究已由传统 A* 或 Hybrid A* 的直接搜索，转向“前端搜索，后端平滑或优化”的两阶段框架\cite{yin2025avp,wu2025agvparking,zhao2024automatedparkingbev}。例如，尹禹兮将改进 Hybrid A* 与后端优化结合，用于提升复杂泊车场景下的路径质量；吴宇轩在改进 A* 基础上引入转弯惩罚和路径跟踪控制；Zhao 则在初始搜索之后继续结合贝塞尔曲线优化和数值优化生成最终轨迹\cite{yin2025avp,wu2025agvparking,zhao2024automatedparkingbev}。研究重点不再只是可达性，研究者开始同时关注轨迹平滑性、执行可行性和狭小空间中的实时性。

但现有泊车文献对场景刻画仍然偏宽。多数研究将泊车统一归入自动泊车或停车场调度问题。研究重点常落在垂直泊车、一般泊车或停车场 AGV 运行\cite{yin2025avp,wu2025agvparking,zhao2024automatedparkingbev}。相较之下，侧方泊车中的特殊几何约束尚未被单独建模。该场景对倒车轨迹形态、边界贴合程度和转向切换更敏感。现有研究对这一细分场景的支撑仍然不足。

\subsection{贝塞尔曲线优化研究现状}

贝塞尔曲线在车辆路径规划中的应用，起初主要承担折线路径平滑任务。已有研究表明，二阶、三阶及五阶贝塞尔曲线均可改善路径连续性，减弱转折突变，并提高轨迹可执行性\cite{zhu2025astarbezier,jiang2019dynamicmotionplanningurban,zhou2025quinticbezier}。其中，朱文亮等将改进 A* 与三阶贝塞尔曲线结合，用于兼顾搜索效率和路径平滑；Jiang 等在城市场景中采用五阶贝塞尔曲线生成满足曲率连续性的参考轨迹；周洋等则从折线路径平滑问题出发，研究了分段五阶贝塞尔曲线及其约束求解方式\cite{zhu2025astarbezier,jiang2019dynamicmotionplanningurban,zhou2025quinticbezier}。低阶贝塞尔曲线便于局部修正，高阶或分段贝塞尔曲线更适合保证曲率连续和整体平滑。

进一步的研究开始把贝塞尔曲线由几何平滑工具，转变为带约束的轨迹优化表达形式。相关文献不再只讨论折线转曲线。研究者开始关注控制点如何满足障碍物、曲率、安全边界和连续时间约束\cite{zhou2025quinticbezier,deolasee2022spatiotemporalbezier,jiang2019dynamicmotionplanningurban}。例如，Deolasee 等在时空走廊内使用分段贝塞尔曲线建模轨迹，并利用其凸包性质保证连续时间安全性；Jiang 等把贝塞尔曲线放入更完整的动态规划框架中，使其同时服务于平滑性与可跟踪性\cite{deolasee2022spatiotemporalbezier,jiang2019dynamicmotionplanningurban}。这说明贝塞尔曲线已由后处理插值，扩展为轨迹优化中的参数化对象。

不过，现有研究多在一般道路、机器人路径平滑或时空联合规划中讨论贝塞尔曲线优点\cite{zhu2025astarbezier,zhou2025quinticbezier,deolasee2022spatiotemporalbezier}。对侧方泊车中的参数化优化问题，现有讨论明显不足。例如，侧方泊车宜采用何种阶数，控制点如何布置，如何同时满足障碍约束与转向约束，这些问题尚缺少专门研究。已有研究证明了贝塞尔曲线具有平滑和优化潜力，但尚未充分回答其在侧方泊车中的最优使用方式。

\subsection{搜索与优化联合规划研究现状}

随着车辆路径规划问题复杂度上升，研究逐渐由“搜索后独立平滑”的串联处理，转向“搜索与优化协同”的联合规划框架。典型做法是先由 A*、Hybrid A* 或其变种生成可行粗路径，再由贝塞尔曲线、时空走廊、数值优化或控制方法对路径进行细化\cite{yin2025avp,li2025hybridastarjoint,deolasee2022spatiotemporalbezier,zhao2024automatedparkingbev}。例如，李妍婷等先用改进混合 A* 搜索初始轨迹，再构造时空走廊并进行后端平滑优化；Deolasee 等在三维时空范围内构建梯形棱柱走廊，并结合分段贝塞尔曲线优化轨迹；Zhao 则通过改进 A* 与后端优化协同提升自动泊车质量\cite{li2025hybridastarjoint,deolasee2022spatiotemporalbezier,zhao2024automatedparkingbev}。这种思路兼顾搜索阶段的全局可达性，也兼顾优化阶段的平滑性和安全性，已成为近年来的重要趋势。

从演化趋势看，这一方向主要有两点变化。其一，前端搜索不再只追求最短路径。研究者开始通过运动学约束、启发式改进和代价函数设计，为后端优化提供更高质量的初始解\cite{zhao2024automatedparkingbev,schichler2024smoothastar,li2025hybridastarjoint}。例如，Schichler 等通过曲率约束和运动原语引导 A* 搜索，Zhao 通过车辆运动学和启发函数优化改善自动泊车初始路径质量\cite{schichler2024smoothastar,zhao2024automatedparkingbev}。其二，后端优化不再局限于局部曲线拟合。研究逐步把时空安全走廊、连续时间约束和控制执行模型纳入统一框架，体现出从“几何平滑”向“约束优化”的明显转变\cite{deolasee2022spatiotemporalbezier,li2025hybridastarjoint}。

然而，这类联合规划框架更多服务于城市动态场景、一般自动泊车或通用低速导航\cite{yin2025avp,li2025hybridastarjoint,zhao2024automatedparkingbev}。针对侧方泊车低速倒车过程的约束刻画仍显不足。现有框架往往强调动态障碍规避、时空耦合或一般停车任务。对边界贴近、倒车轨迹变形和有限空间转向切换的讨论较少。搜索与优化联合规划为你的研究提供了总体技术路线，但在侧方泊车这一具体问题上仍有明显细化空间。

\subsection{小结}

综合已有研究可以看出，现有文献已形成三点共识。其一，A* 与 Hybrid A* 适合作为车辆泊车路径的初始搜索工具\cite{yin2025avp,wu2025agvparking,zhao2024automatedparkingbev}。其二，贝塞尔曲线适合改善路径平滑性、曲率连续性和可执行性\cite{zhu2025astarbezier,zhou2025quinticbezier,deolasee2022spatiotemporalbezier}。其三，搜索与优化耦合已成为提升规划质量的重要框架\cite{li2025hybridastarjoint,deolasee2022spatiotemporalbezier,zhao2024automatedparkingbev}。

但现有研究仍存在三个直接空白。第一，侧方泊车场景本身缺少专门研究。第二，贝塞尔曲线在侧方泊车中的参数化优化问题尚未充分展开。第三，现有搜索--优化联合框架对低速倒车和狭小空间边界约束的适配不足。基于上述不足，围绕侧方泊车中的贝塞尔曲线优化开展研究，既能承接现有技术路线，也具有明确的问题针对性。
 引入本文件。
% 2. 参考文献数据库建议加入 references/related_work_gap123.bib。

\chapter{相关研究综述}

侧方泊车是低速自动驾驶中的典型任务。该任务对路径规划提出了较高要求。车辆既要在有限空间内完成无碰撞机动，又要保证轨迹平滑、转向可执行和末端停放精度。传统搜索算法易于生成可行路径，但所得结果常呈折线形态，难以直接满足车辆连续转向需求。贝塞尔曲线具有良好的平滑性与参数化表达能力，因而被广泛用于车辆轨迹优化。围绕侧方泊车中的贝塞尔曲线优化问题，梳理现有研究进展，有助于明确相关方法的演化脉络，并为后续问题建模与算法设计提供依据。

\subsection{侧方泊车相关研究现状}

现有泊车路径规划研究，主要面向自动代客泊车与停车场低速机动场景\cite{yin2025avp,wu2025agvparking,zhao2024automatedparkingbev}。相关工作多以环境建模、初始路径搜索和轨迹跟踪控制为主线展开。研究目标已从静态最短路搜索，转向考虑车辆运动学、障碍分布和控制执行约束的综合规划。其中，尹禹兮围绕自主代客泊车构建了搜索与优化结合的分层规划思路，吴宇轩则以智能停车场 AGV 为对象，讨论了路径规划与路径跟踪的一体化实现，Zhao 进一步把视觉 BEV 局部建图引入自动泊车规划过程\cite{yin2025avp,wu2025agvparking,zhao2024automatedparkingbev}。

从方法演化看，泊车研究已由传统 A* 或 Hybrid A* 的直接搜索，转向“前端搜索，后端平滑或优化”的两阶段框架\cite{yin2025avp,wu2025agvparking,zhao2024automatedparkingbev}。例如，尹禹兮将改进 Hybrid A* 与后端优化结合，用于提升复杂泊车场景下的路径质量；吴宇轩在改进 A* 基础上引入转弯惩罚和路径跟踪控制；Zhao 则在初始搜索之后继续结合贝塞尔曲线优化和数值优化生成最终轨迹\cite{yin2025avp,wu2025agvparking,zhao2024automatedparkingbev}。研究重点不再只是可达性，研究者开始同时关注轨迹平滑性、执行可行性和狭小空间中的实时性。

但现有泊车文献对场景刻画仍然偏宽。多数研究将泊车统一归入自动泊车或停车场调度问题。研究重点常落在垂直泊车、一般泊车或停车场 AGV 运行\cite{yin2025avp,wu2025agvparking,zhao2024automatedparkingbev}。相较之下，侧方泊车中的特殊几何约束尚未被单独建模。该场景对倒车轨迹形态、边界贴合程度和转向切换更敏感。现有研究对这一细分场景的支撑仍然不足。

\subsection{贝塞尔曲线优化研究现状}

贝塞尔曲线在车辆路径规划中的应用，起初主要承担折线路径平滑任务。已有研究表明，二阶、三阶及五阶贝塞尔曲线均可改善路径连续性，减弱转折突变，并提高轨迹可执行性\cite{zhu2025astarbezier,jiang2019dynamicmotionplanningurban,zhou2025quinticbezier}。其中，朱文亮等将改进 A* 与三阶贝塞尔曲线结合，用于兼顾搜索效率和路径平滑；Jiang 等在城市场景中采用五阶贝塞尔曲线生成满足曲率连续性的参考轨迹；周洋等则从折线路径平滑问题出发，研究了分段五阶贝塞尔曲线及其约束求解方式\cite{zhu2025astarbezier,jiang2019dynamicmotionplanningurban,zhou2025quinticbezier}。低阶贝塞尔曲线便于局部修正，高阶或分段贝塞尔曲线更适合保证曲率连续和整体平滑。

进一步的研究开始把贝塞尔曲线由几何平滑工具，转变为带约束的轨迹优化表达形式。相关文献不再只讨论折线转曲线。研究者开始关注控制点如何满足障碍物、曲率、安全边界和连续时间约束\cite{zhou2025quinticbezier,deolasee2022spatiotemporalbezier,jiang2019dynamicmotionplanningurban}。例如，Deolasee 等在时空走廊内使用分段贝塞尔曲线建模轨迹，并利用其凸包性质保证连续时间安全性；Jiang 等把贝塞尔曲线放入更完整的动态规划框架中，使其同时服务于平滑性与可跟踪性\cite{deolasee2022spatiotemporalbezier,jiang2019dynamicmotionplanningurban}。这说明贝塞尔曲线已由后处理插值，扩展为轨迹优化中的参数化对象。

不过，现有研究多在一般道路、机器人路径平滑或时空联合规划中讨论贝塞尔曲线优点\cite{zhu2025astarbezier,zhou2025quinticbezier,deolasee2022spatiotemporalbezier}。对侧方泊车中的参数化优化问题，现有讨论明显不足。例如，侧方泊车宜采用何种阶数，控制点如何布置，如何同时满足障碍约束与转向约束，这些问题尚缺少专门研究。已有研究证明了贝塞尔曲线具有平滑和优化潜力，但尚未充分回答其在侧方泊车中的最优使用方式。

\subsection{搜索与优化联合规划研究现状}

随着车辆路径规划问题复杂度上升，研究逐渐由“搜索后独立平滑”的串联处理，转向“搜索与优化协同”的联合规划框架。典型做法是先由 A*、Hybrid A* 或其变种生成可行粗路径，再由贝塞尔曲线、时空走廊、数值优化或控制方法对路径进行细化\cite{yin2025avp,li2025hybridastarjoint,deolasee2022spatiotemporalbezier,zhao2024automatedparkingbev}。例如，李妍婷等先用改进混合 A* 搜索初始轨迹，再构造时空走廊并进行后端平滑优化；Deolasee 等在三维时空范围内构建梯形棱柱走廊，并结合分段贝塞尔曲线优化轨迹；Zhao 则通过改进 A* 与后端优化协同提升自动泊车质量\cite{li2025hybridastarjoint,deolasee2022spatiotemporalbezier,zhao2024automatedparkingbev}。这种思路兼顾搜索阶段的全局可达性，也兼顾优化阶段的平滑性和安全性，已成为近年来的重要趋势。

从演化趋势看，这一方向主要有两点变化。其一，前端搜索不再只追求最短路径。研究者开始通过运动学约束、启发式改进和代价函数设计，为后端优化提供更高质量的初始解\cite{zhao2024automatedparkingbev,schichler2024smoothastar,li2025hybridastarjoint}。例如，Schichler 等通过曲率约束和运动原语引导 A* 搜索，Zhao 通过车辆运动学和启发函数优化改善自动泊车初始路径质量\cite{schichler2024smoothastar,zhao2024automatedparkingbev}。其二，后端优化不再局限于局部曲线拟合。研究逐步把时空安全走廊、连续时间约束和控制执行模型纳入统一框架，体现出从“几何平滑”向“约束优化”的明显转变\cite{deolasee2022spatiotemporalbezier,li2025hybridastarjoint}。

然而，这类联合规划框架更多服务于城市动态场景、一般自动泊车或通用低速导航\cite{yin2025avp,li2025hybridastarjoint,zhao2024automatedparkingbev}。针对侧方泊车低速倒车过程的约束刻画仍显不足。现有框架往往强调动态障碍规避、时空耦合或一般停车任务。对边界贴近、倒车轨迹变形和有限空间转向切换的讨论较少。搜索与优化联合规划为你的研究提供了总体技术路线，但在侧方泊车这一具体问题上仍有明显细化空间。

\subsection{小结}

综合已有研究可以看出，现有文献已形成三点共识。其一，A* 与 Hybrid A* 适合作为车辆泊车路径的初始搜索工具\cite{yin2025avp,wu2025agvparking,zhao2024automatedparkingbev}。其二，贝塞尔曲线适合改善路径平滑性、曲率连续性和可执行性\cite{zhu2025astarbezier,zhou2025quinticbezier,deolasee2022spatiotemporalbezier}。其三，搜索与优化耦合已成为提升规划质量的重要框架\cite{li2025hybridastarjoint,deolasee2022spatiotemporalbezier,zhao2024automatedparkingbev}。

但现有研究仍存在三个直接空白。第一，侧方泊车场景本身缺少专门研究。第二，贝塞尔曲线在侧方泊车中的参数化优化问题尚未充分展开。第三，现有搜索--优化联合框架对低速倒车和狭小空间边界约束的适配不足。基于上述不足，围绕侧方泊车中的贝塞尔曲线优化开展研究，既能承接现有技术路线，也具有明确的问题针对性。
 引入本文件。
% 2. 参考文献数据库建议加入 references/related_work_gap123.bib。

\chapter{相关研究综述}

侧方泊车是低速自动驾驶中的典型任务。该任务对路径规划提出了较高要求。车辆既要在有限空间内完成无碰撞机动，又要保证轨迹平滑、转向可执行和末端停放精度。传统搜索算法易于生成可行路径，但所得结果常呈折线形态，难以直接满足车辆连续转向需求。贝塞尔曲线具有良好的平滑性与参数化表达能力，因而被广泛用于车辆轨迹优化。围绕侧方泊车中的贝塞尔曲线优化问题，梳理现有研究进展，有助于明确相关方法的演化脉络，并为后续问题建模与算法设计提供依据。

\subsection{侧方泊车相关研究现状}

现有泊车路径规划研究，主要面向自动代客泊车与停车场低速机动场景\cite{yin2025avp,wu2025agvparking,zhao2024automatedparkingbev}。相关工作多以环境建模、初始路径搜索和轨迹跟踪控制为主线展开。研究目标已从静态最短路搜索，转向考虑车辆运动学、障碍分布和控制执行约束的综合规划。其中，尹禹兮围绕自主代客泊车构建了搜索与优化结合的分层规划思路，吴宇轩则以智能停车场 AGV 为对象，讨论了路径规划与路径跟踪的一体化实现，Zhao 进一步把视觉 BEV 局部建图引入自动泊车规划过程\cite{yin2025avp,wu2025agvparking,zhao2024automatedparkingbev}。

从方法演化看，泊车研究已由传统 A* 或 Hybrid A* 的直接搜索，转向“前端搜索，后端平滑或优化”的两阶段框架\cite{yin2025avp,wu2025agvparking,zhao2024automatedparkingbev}。例如，尹禹兮将改进 Hybrid A* 与后端优化结合，用于提升复杂泊车场景下的路径质量；吴宇轩在改进 A* 基础上引入转弯惩罚和路径跟踪控制；Zhao 则在初始搜索之后继续结合贝塞尔曲线优化和数值优化生成最终轨迹\cite{yin2025avp,wu2025agvparking,zhao2024automatedparkingbev}。研究重点不再只是可达性，研究者开始同时关注轨迹平滑性、执行可行性和狭小空间中的实时性。

但现有泊车文献对场景刻画仍然偏宽。多数研究将泊车统一归入自动泊车或停车场调度问题。研究重点常落在垂直泊车、一般泊车或停车场 AGV 运行\cite{yin2025avp,wu2025agvparking,zhao2024automatedparkingbev}。相较之下，侧方泊车中的特殊几何约束尚未被单独建模。该场景对倒车轨迹形态、边界贴合程度和转向切换更敏感。现有研究对这一细分场景的支撑仍然不足。

\subsection{贝塞尔曲线优化研究现状}

贝塞尔曲线在车辆路径规划中的应用，起初主要承担折线路径平滑任务。已有研究表明，二阶、三阶及五阶贝塞尔曲线均可改善路径连续性，减弱转折突变，并提高轨迹可执行性\cite{zhu2025astarbezier,jiang2019dynamicmotionplanningurban,zhou2025quinticbezier}。其中，朱文亮等将改进 A* 与三阶贝塞尔曲线结合，用于兼顾搜索效率和路径平滑；Jiang 等在城市场景中采用五阶贝塞尔曲线生成满足曲率连续性的参考轨迹；周洋等则从折线路径平滑问题出发，研究了分段五阶贝塞尔曲线及其约束求解方式\cite{zhu2025astarbezier,jiang2019dynamicmotionplanningurban,zhou2025quinticbezier}。低阶贝塞尔曲线便于局部修正，高阶或分段贝塞尔曲线更适合保证曲率连续和整体平滑。

进一步的研究开始把贝塞尔曲线由几何平滑工具，转变为带约束的轨迹优化表达形式。相关文献不再只讨论折线转曲线。研究者开始关注控制点如何满足障碍物、曲率、安全边界和连续时间约束\cite{zhou2025quinticbezier,deolasee2022spatiotemporalbezier,jiang2019dynamicmotionplanningurban}。例如，Deolasee 等在时空走廊内使用分段贝塞尔曲线建模轨迹，并利用其凸包性质保证连续时间安全性；Jiang 等把贝塞尔曲线放入更完整的动态规划框架中，使其同时服务于平滑性与可跟踪性\cite{deolasee2022spatiotemporalbezier,jiang2019dynamicmotionplanningurban}。这说明贝塞尔曲线已由后处理插值，扩展为轨迹优化中的参数化对象。

不过，现有研究多在一般道路、机器人路径平滑或时空联合规划中讨论贝塞尔曲线优点\cite{zhu2025astarbezier,zhou2025quinticbezier,deolasee2022spatiotemporalbezier}。对侧方泊车中的参数化优化问题，现有讨论明显不足。例如，侧方泊车宜采用何种阶数，控制点如何布置，如何同时满足障碍约束与转向约束，这些问题尚缺少专门研究。已有研究证明了贝塞尔曲线具有平滑和优化潜力，但尚未充分回答其在侧方泊车中的最优使用方式。

\subsection{搜索与优化联合规划研究现状}

随着车辆路径规划问题复杂度上升，研究逐渐由“搜索后独立平滑”的串联处理，转向“搜索与优化协同”的联合规划框架。典型做法是先由 A*、Hybrid A* 或其变种生成可行粗路径，再由贝塞尔曲线、时空走廊、数值优化或控制方法对路径进行细化\cite{yin2025avp,li2025hybridastarjoint,deolasee2022spatiotemporalbezier,zhao2024automatedparkingbev}。例如，李妍婷等先用改进混合 A* 搜索初始轨迹，再构造时空走廊并进行后端平滑优化；Deolasee 等在三维时空范围内构建梯形棱柱走廊，并结合分段贝塞尔曲线优化轨迹；Zhao 则通过改进 A* 与后端优化协同提升自动泊车质量\cite{li2025hybridastarjoint,deolasee2022spatiotemporalbezier,zhao2024automatedparkingbev}。这种思路兼顾搜索阶段的全局可达性，也兼顾优化阶段的平滑性和安全性，已成为近年来的重要趋势。

从演化趋势看，这一方向主要有两点变化。其一，前端搜索不再只追求最短路径。研究者开始通过运动学约束、启发式改进和代价函数设计，为后端优化提供更高质量的初始解\cite{zhao2024automatedparkingbev,schichler2024smoothastar,li2025hybridastarjoint}。例如，Schichler 等通过曲率约束和运动原语引导 A* 搜索，Zhao 通过车辆运动学和启发函数优化改善自动泊车初始路径质量\cite{schichler2024smoothastar,zhao2024automatedparkingbev}。其二，后端优化不再局限于局部曲线拟合。研究逐步把时空安全走廊、连续时间约束和控制执行模型纳入统一框架，体现出从“几何平滑”向“约束优化”的明显转变\cite{deolasee2022spatiotemporalbezier,li2025hybridastarjoint}。

然而，这类联合规划框架更多服务于城市动态场景、一般自动泊车或通用低速导航\cite{yin2025avp,li2025hybridastarjoint,zhao2024automatedparkingbev}。针对侧方泊车低速倒车过程的约束刻画仍显不足。现有框架往往强调动态障碍规避、时空耦合或一般停车任务。对边界贴近、倒车轨迹变形和有限空间转向切换的讨论较少。搜索与优化联合规划为你的研究提供了总体技术路线，但在侧方泊车这一具体问题上仍有明显细化空间。

\subsection{小结}

综合已有研究可以看出，现有文献已形成三点共识。其一，A* 与 Hybrid A* 适合作为车辆泊车路径的初始搜索工具\cite{yin2025avp,wu2025agvparking,zhao2024automatedparkingbev}。其二，贝塞尔曲线适合改善路径平滑性、曲率连续性和可执行性\cite{zhu2025astarbezier,zhou2025quinticbezier,deolasee2022spatiotemporalbezier}。其三，搜索与优化耦合已成为提升规划质量的重要框架\cite{li2025hybridastarjoint,deolasee2022spatiotemporalbezier,zhao2024automatedparkingbev}。

但现有研究仍存在三个直接空白。第一，侧方泊车场景本身缺少专门研究。第二，贝塞尔曲线在侧方泊车中的参数化优化问题尚未充分展开。第三，现有搜索--优化联合框架对低速倒车和狭小空间边界约束的适配不足。基于上述不足，围绕侧方泊车中的贝塞尔曲线优化开展研究，既能承接现有技术路线，也具有明确的问题针对性。
 引入本文件。
% 2. 参考文献数据库建议加入 references/related_work_gap123.bib。

\chapter{相关研究综述}

侧方泊车是低速自动驾驶中的典型任务。该任务对路径规划提出了较高要求。车辆既要在有限空间内完成无碰撞机动，又要保证轨迹平滑、转向可执行和末端停放精度。传统搜索算法易于生成可行路径，但所得结果常呈折线形态，难以直接满足车辆连续转向需求。贝塞尔曲线具有良好的平滑性与参数化表达能力，因而被广泛用于车辆轨迹优化。围绕侧方泊车中的贝塞尔曲线优化问题，梳理现有研究进展，有助于明确相关方法的演化脉络，并为后续问题建模与算法设计提供依据。

\subsection{侧方泊车相关研究现状}

现有泊车路径规划研究，主要面向自动代客泊车与停车场低速机动场景\cite{yin2025avp,wu2025agvparking,zhao2024automatedparkingbev}。相关工作多以环境建模、初始路径搜索和轨迹跟踪控制为主线展开。研究目标已从静态最短路搜索，转向考虑车辆运动学、障碍分布和控制执行约束的综合规划。其中，尹禹兮围绕自主代客泊车构建了搜索与优化结合的分层规划思路，吴宇轩则以智能停车场 AGV 为对象，讨论了路径规划与路径跟踪的一体化实现，Zhao 进一步把视觉 BEV 局部建图引入自动泊车规划过程\cite{yin2025avp,wu2025agvparking,zhao2024automatedparkingbev}。

从方法演化看，泊车研究已由传统 A* 或 Hybrid A* 的直接搜索，转向“前端搜索，后端平滑或优化”的两阶段框架\cite{yin2025avp,wu2025agvparking,zhao2024automatedparkingbev}。例如，尹禹兮将改进 Hybrid A* 与后端优化结合，用于提升复杂泊车场景下的路径质量；吴宇轩在改进 A* 基础上引入转弯惩罚和路径跟踪控制；Zhao 则在初始搜索之后继续结合贝塞尔曲线优化和数值优化生成最终轨迹\cite{yin2025avp,wu2025agvparking,zhao2024automatedparkingbev}。研究重点不再只是可达性，研究者开始同时关注轨迹平滑性、执行可行性和狭小空间中的实时性。

但现有泊车文献对场景刻画仍然偏宽。多数研究将泊车统一归入自动泊车或停车场调度问题。研究重点常落在垂直泊车、一般泊车或停车场 AGV 运行\cite{yin2025avp,wu2025agvparking,zhao2024automatedparkingbev}。相较之下，侧方泊车中的特殊几何约束尚未被单独建模。该场景对倒车轨迹形态、边界贴合程度和转向切换更敏感。现有研究对这一细分场景的支撑仍然不足。

\subsection{贝塞尔曲线优化研究现状}

贝塞尔曲线在车辆路径规划中的应用，起初主要承担折线路径平滑任务。已有研究表明，二阶、三阶及五阶贝塞尔曲线均可改善路径连续性，减弱转折突变，并提高轨迹可执行性\cite{zhu2025astarbezier,jiang2019dynamicmotionplanningurban,zhou2025quinticbezier}。其中，朱文亮等将改进 A* 与三阶贝塞尔曲线结合，用于兼顾搜索效率和路径平滑；Jiang 等在城市场景中采用五阶贝塞尔曲线生成满足曲率连续性的参考轨迹；周洋等则从折线路径平滑问题出发，研究了分段五阶贝塞尔曲线及其约束求解方式\cite{zhu2025astarbezier,jiang2019dynamicmotionplanningurban,zhou2025quinticbezier}。低阶贝塞尔曲线便于局部修正，高阶或分段贝塞尔曲线更适合保证曲率连续和整体平滑。

进一步的研究开始把贝塞尔曲线由几何平滑工具，转变为带约束的轨迹优化表达形式。相关文献不再只讨论折线转曲线。研究者开始关注控制点如何满足障碍物、曲率、安全边界和连续时间约束\cite{zhou2025quinticbezier,deolasee2022spatiotemporalbezier,jiang2019dynamicmotionplanningurban}。例如，Deolasee 等在时空走廊内使用分段贝塞尔曲线建模轨迹，并利用其凸包性质保证连续时间安全性；Jiang 等把贝塞尔曲线放入更完整的动态规划框架中，使其同时服务于平滑性与可跟踪性\cite{deolasee2022spatiotemporalbezier,jiang2019dynamicmotionplanningurban}。这说明贝塞尔曲线已由后处理插值，扩展为轨迹优化中的参数化对象。

不过，现有研究多在一般道路、机器人路径平滑或时空联合规划中讨论贝塞尔曲线优点\cite{zhu2025astarbezier,zhou2025quinticbezier,deolasee2022spatiotemporalbezier}。对侧方泊车中的参数化优化问题，现有讨论明显不足。例如，侧方泊车宜采用何种阶数，控制点如何布置，如何同时满足障碍约束与转向约束，这些问题尚缺少专门研究。已有研究证明了贝塞尔曲线具有平滑和优化潜力，但尚未充分回答其在侧方泊车中的最优使用方式。

\subsection{搜索与优化联合规划研究现状}

随着车辆路径规划问题复杂度上升，研究逐渐由“搜索后独立平滑”的串联处理，转向“搜索与优化协同”的联合规划框架。典型做法是先由 A*、Hybrid A* 或其变种生成可行粗路径，再由贝塞尔曲线、时空走廊、数值优化或控制方法对路径进行细化\cite{yin2025avp,li2025hybridastarjoint,deolasee2022spatiotemporalbezier,zhao2024automatedparkingbev}。例如，李妍婷等先用改进混合 A* 搜索初始轨迹，再构造时空走廊并进行后端平滑优化；Deolasee 等在三维时空范围内构建梯形棱柱走廊，并结合分段贝塞尔曲线优化轨迹；Zhao 则通过改进 A* 与后端优化协同提升自动泊车质量\cite{li2025hybridastarjoint,deolasee2022spatiotemporalbezier,zhao2024automatedparkingbev}。这种思路兼顾搜索阶段的全局可达性，也兼顾优化阶段的平滑性和安全性，已成为近年来的重要趋势。

从演化趋势看，这一方向主要有两点变化。其一，前端搜索不再只追求最短路径。研究者开始通过运动学约束、启发式改进和代价函数设计，为后端优化提供更高质量的初始解\cite{zhao2024automatedparkingbev,schichler2024smoothastar,li2025hybridastarjoint}。例如，Schichler 等通过曲率约束和运动原语引导 A* 搜索，Zhao 通过车辆运动学和启发函数优化改善自动泊车初始路径质量\cite{schichler2024smoothastar,zhao2024automatedparkingbev}。其二，后端优化不再局限于局部曲线拟合。研究逐步把时空安全走廊、连续时间约束和控制执行模型纳入统一框架，体现出从“几何平滑”向“约束优化”的明显转变\cite{deolasee2022spatiotemporalbezier,li2025hybridastarjoint}。

然而，这类联合规划框架更多服务于城市动态场景、一般自动泊车或通用低速导航\cite{yin2025avp,li2025hybridastarjoint,zhao2024automatedparkingbev}。针对侧方泊车低速倒车过程的约束刻画仍显不足。现有框架往往强调动态障碍规避、时空耦合或一般停车任务。对边界贴近、倒车轨迹变形和有限空间转向切换的讨论较少。搜索与优化联合规划为你的研究提供了总体技术路线，但在侧方泊车这一具体问题上仍有明显细化空间。

\subsection{小结}

综合已有研究可以看出，现有文献已形成三点共识。其一，A* 与 Hybrid A* 适合作为车辆泊车路径的初始搜索工具\cite{yin2025avp,wu2025agvparking,zhao2024automatedparkingbev}。其二，贝塞尔曲线适合改善路径平滑性、曲率连续性和可执行性\cite{zhu2025astarbezier,zhou2025quinticbezier,deolasee2022spatiotemporalbezier}。其三，搜索与优化耦合已成为提升规划质量的重要框架\cite{li2025hybridastarjoint,deolasee2022spatiotemporalbezier,zhao2024automatedparkingbev}。

但现有研究仍存在三个直接空白。第一，侧方泊车场景本身缺少专门研究。第二，贝塞尔曲线在侧方泊车中的参数化优化问题尚未充分展开。第三，现有搜索--优化联合框架对低速倒车和狭小空间边界约束的适配不足。基于上述不足，围绕侧方泊车中的贝塞尔曲线优化开展研究，既能承接现有技术路线，也具有明确的问题针对性。
